\paragraph{The integral light lattice in the marked chart.}
\label{par:slice-light-integral}
Fix a lift and either of the two runs, with construction field
\(K_0\), half-level ideal \(\mathfrak m_b\), and full level
\(M=g_0g_1\). Write \(a_{\tau,i}\) for the entries of a weight
\(w\) on the \(\tau\)-eigenspace and \(d_\tau=\sum_i a_{\tau,i}\).
Use \(\mathcal T_w\) for the theta rules with these entries and
the prescribed number of copies. A rule uses an integral harmonic
polynomial \(P\), of the specified flag weight, and marked level
points on the copies.

For each marked point take an \(M\)-division
point killed by \(\mathfrak m_b^2\). The normalized multiplication
isomorphism and the symmetric theta lift give a translated section
\(H_b\) of the \(M^2\)-th power of the rigidified theta line on
copy \(b\). The rule reads
\begin{equation}\label{eq:light-rule-value}
 M^{-\deg P}\sum_{\boldsymbol\mu} P_{\boldsymbol\mu}
        \prod_b \nabla^{\boldsymbol\mu_b} H_b(0).
\end{equation}
Here \(\boldsymbol\mu_b\) uses the weight-labelled tangent
directions on copy \(b\), and \(\nabla\) is the covariant
derivative in the given frame. The translation and trivialization
laws give this expression after a faithfully flat cover carrying
the marked sections. The harmonic and theta transformation laws
identify it with the extension of the original rule value for
any admissible choice of the \(H_b\). The power of \(M\) in
\eqref{eq:light-rule-value} is part of the rule throughout the
comparison. \label{cor:slice-light-integral-1}

\begin{proposition}[Integral evaluation of the scaled lattice of theta rules]
\label{prop:scaled-integral-light-reading}
After scalar multiplication by \(a_W\) in \eqref{eq:light-global-scalar},
the chosen light family generates the image of the integral span
of \(\mathcal T_{w_L}\). Let \(w_l\) be the scaled light generators and \(b_T\) the scaled value of a rule \(T\). At every allowed integral framed point \(y\), in the fixed orientation and coefficient embedding,
\[
 \min_l\operatorname{ord}w_l(y)=\inf_{T\in\mathcal T_{w_L}}\operatorname{ord}b_T(y).
\]

\end{proposition}

\begin{proof}
\proofstep{Integral frames and division by the level scalar.}
Take an algebraic integral framed point \(y\) over a complete
DVR \(V\) above \(\lambda\), with relative ramification \(e\).
Its residue field is finite. Choose one of its integral models,
with polarized semiabelian group \(G\). The split base-change
theorem for the semiabelian model over this complete DVR gives
a finite local extension \(V'/V\) with finite residue field and a
split uniformization over \(V'\). The integral-model comparison
transports that uniformization and its torsion identification to
every integral model of the extended point.

Choose the
integral formal-group chart at the pulled-back frame. Its
polarization generators are compatible under jet truncation
over \(V'\).
For example, for jet orders \(a\le b\), the compatible reductions of the chosen formal parameters give
\[
  V'[\![X]\!]/(X)^{b+1}\longrightarrow V'[\![X]\!]/(X)^{a+1}.
\]
 The framed base-change map pulls back the
underlying point, all the labelled tangent vectors and the
theta-line generator. The incidence map respects this frame
extension, so the new point lies over the same point in the
boundary comparison. \label{cor:slice-light-integral-2}

Let \(\iota:\operatorname{Frac}V\hookrightarrow\operatorname{Frac}V'\)
be the extension. The normalized reading commutes with \(\iota\).
Indeed, for a rule at the new frame choose a field over
\(\operatorname{Frac}V'\) carrying the translated sections. Such
a field is obtained by taking a residue field at a maximal ideal of
the faithfully flat coefficient algebra and pulling back its
sections. Over this field, the composite frame map from the old
point and the frame map from the new point give the same
expression \eqref{eq:light-rule-value} using the same sections.
Injectivity of the field map proves

\[
  f(y')=\iota(f(y)).
\]
 This works for each rule and each weight
on the same \(V'\). The auxiliary field used to test the equality
can vary with the rule. If \(r\) is the ramification of
\(V'/V\), then
\begin{equation}\label{eq:light-frame-order}
 \frac{\operatorname{ord}_{V'} f(y')}{re}
        =\frac{\operatorname{ord}_{V} f(y)}{e}.
\end{equation}
Both sides are infinite when the reading is zero. All the
integral germs and the local interpolation estimates below may
thus be taken on the split chart, with its own ramification.
\label{cor:slice-light-integral-3}

\proofstep{Using the half-ideal of the selected orientation.}
For the actual orientation choose the elements \(\pi_i\) of
the fixed imaginary quadratic subfield used by the run, with
\(\pi_i\bar\pi_i=g_i\), and put \(\alpha=\pi_0\pi_1\).
The oriented markings give
\(\mathfrak m_b=(\bar\alpha)\) and 
\[
  \alpha\bar\alpha=M.
\]
In fact each \(\bar\pi_i\) divides a generator of the half
level. Their ideals are coprime because \(g_0,g_1\) are distinct,
so their product divides that generator. The relative norm of the
quotient is one, since the generator and \(\bar\alpha\) have the
same relative norm \(M\). Thus the quotient is a unit. Away
from the level one can equivalently use any generator with this
relative norm.

At the level use \(\alpha\) with the indicated
orientation. Let \(\iota_0\) denote the coefficient embedding of
the fixed imaginary quadratic field defining the plus type.
Each plus embedding agrees with \(\iota_0\) there and each
minus embedding is its conjugate. Put
\(D^\pm(w)=\sum_{\tau\text{ of type }\pm}d_\tau\). The weight
determinant of the action of \(\alpha\) on the labelled frame is
\begin{align}
 \mathcal D_w(\alpha)
     &=\prod_\tau\tau_K(\alpha)^{d_\tau} \notag\\
     &=\prod_{i=0}^1
         \iota_0(\pi_i)^{D^+(w)}\,
         \iota_0(\bar\pi_i)^{D^-(w)},
   \label{eq:light-oriented-determinant}
\end{align}
where \(\tau_K\) are the embeddings used in the coefficient
field. The first formula is the action on each row degree of
\(P\). The second follows by summing those degrees in the two
signature blocks. In particular \(D_F=D^+(w_L)+D^-(w_L)\).
\label{cor:slice-light-integral-4}

\proofstep{The minimum valuation at the smaller weight.}
At the smaller weight, the scalar multiplying the lattice is
\begin{equation}\label{eq:light-global-scalar}
 a_W=\mathcal D_{w_L}(\alpha)\,
       Q_S^{D_F}q^{\lceil D_F/2\rceil}d^{\lceil D_F/2\rceil}.
\end{equation}
The local value of \(a_W\) is obtained from this coefficient
scalar by the same embedding as the rule. The light lattice is
\(a_W\operatorname{span}_{\mathcal O_{K_0}}\mathcal T_{w_L}\),
using the cusp inclusion at this weight. This equality follows
by taking the integral span of the rules and applying the scalar
map to the span. The chosen family \((w_l)\) generates exactly
this image, and the new choice of short sections retains it.

In the fixed integral frame, an integral combination has
valuation at least the least valuation of its terms. Hence the
least valuation of the scaled evaluations is attained on a generator \(w_l\). Equations \eqref{eq:light-frame-order}
and \eqref{eq:light-oriented-determinant} apply at every allowed
weight, including the two weights in the full and truncated
evaluation comparisons. The choice of chart depends on the
point and marking. It introduces no weight-dependent scaling
into these evaluations. \label{cor:slice-light-integral-5}
\end{proof}
